% ID: SEL-D0039-OB-E02
% ESTADO: validada. S040, cotejada con examen y criterios.
\textbf{a)} Cada tramo es continuo en su parte del dominio.
En $2$ el primero vale 6 y el límite derecho es $16+a$.
Igualarlos da $\boxed{a=-10}$, con continuidad en $[0,\infty)$.

\textbf{b)} Para $x>2$, con $a=-10$ podemos escribir
$f(x)=8-2/(x-1)$. Entonces $f'(x)=2/(x-1)^2$
y $f'(3)=1/2>0$. Sí es creciente en un entorno de 3.

\textbf{c)} El segundo tramo tiende a 8 cuando $x\to+\infty$,
de modo que $\boxed{y=8}$ es asíntota horizontal.
El denominador sólo se anularía en $1$, fuera del tramo racional.
Como la función es continua en todo su dominio y tiene límite
finito en $0$ por la derecha, no hay asíntotas verticales.
